% ============================================================
%  论文主文档 —— Springer / MTAP 版（当前主攻目标）
%  编译： make mtap
%  目标：Multimedia Tools and Applications, Track 10
%        (generative AI in multimedia production pipelines)
%  注意：期刊版章节顺序自由，但需补 Related Work 的横向对比表。
% ============================================================
\documentclass[runningheads]{llncs}

% ============================================================
%  共用宏包与命令 —— 两个 main.tex 都 \input 这个文件
%  作者单位：真实值放在本地 affiliation.tex（已被 .gitignore 拦截）。
%    · 有该文件  → 用真实单位（本机编译）
%    · 无该文件  → 用下面占位值（公开仓库 / 别人 clone 后编译）
% ============================================================

\usepackage[T1]{fontenc}
\usepackage[utf8]{inputenc}
\usepackage{booktabs}       % 三线表
\usepackage{graphicx}
\usepackage{url}
\usepackage{amsmath,amssymb}
\usepackage{booktabs,xcolor,colortbl}
\usepackage{hyperref}
\hypersetup{colorlinks=true,linkcolor=blue,citecolor=blue,urlcolor=blue}

% ---------- 自定义命令（全篇统一，不要在章节里另写） ----------
\newcommand{\agnes}{\textsc{Agnes}}
\newcommand{\comfy}{\textsc{ComfyUI}}
\newcommand{\styledef}{STYLE\_LOCK}
\newcommand{\insight}{InsightFace}
\newcommand{\vds}{VDS}

\newcommand{\nts}{\textbf{S1E01}}

% 行内代码统一风格
\newcommand{\code}[1]{\texttt{#1}}

% 场景代号强调（论文里会反复出现场景号，别写漏下标）
\newcommand{\shot}{\textsc{Shot}}

% ---------- 图表默认参数 ----------
\graphicspath{{figures/}}
\setcounter{topnumber}{3}
\setcounter{bottomnumber}{3}
\setcounter{totalnumber}{6}
\renewcommand{\topfraction}{0.9}
\renewcommand{\bottomfraction}{0.9}

% ---------- 作者 / 单位 ----------
% 真实单位只在本地 affiliation.tex。该文件被 .gitignore 拦截，不会进公开仓库。
% 顺序不能颠倒：先读真值，再用 \providecommand 补占位（已定义的不覆盖）。
\IfFileExists{affiliation.tex}{\input{affiliation.tex}}{}
\providecommand{\theauthor}{[Author name withheld until submission]}
\providecommand{\theaffil}{[Affiliation withheld until submission]}
\providecommand{\theemail}{[email]}

% ---------- 语言与排版 ----------
\hyphenation{ComfyUI InsightFace negative-prompt}


\begin{document}

\title{An Industrial-Grade Generative Pipeline for Character-Consistent
       Animated Series: System Design, Failure Taxonomy, and
       Negative-Prompt Efficacy}

% ---------- 作者 / 单位（Springer 格式，真实值见本地 affiliation.tex） ----------
\author{\theauthor\inst{1}}
\institute{\theaffil}
\email{\theemail}

\maketitle

\begin{abstract}
% TODO: 150-200 词，结构块化（ significance / contributions / results ）
\end{abstract}

\begin{keywords}
Generative AI, animation production pipeline, character consistency,
failure analysis, negative prompt, quality control
\end{keywords}

\section{Introduction}

\subsection{Motivation}

Producing a character-consistent animated episode with generative models
requires sustaining visual identity across dozens of shots, each one
regenerated from scratch by text-to-image and text-to-video models. In our
production of \nts, the observed one-pass acceptance rate was
\textbf{12 of 52 shots}
(\textbf{23\%}). The remaining 40 shots required re-generation, and a
non-trivial fraction were rejected more than once.

This figure is, to the best of our knowledge, \emph{not} unusual. It is,
however, almost never reported. Production teams treat per-shot failure as an
operational cost and have no incentive to publish its distribution. The
result is a field operating without a shared quantitative account of
\emph{what actually breaks}.

\subsection{Contributions}

We contribute a system description, an empirical failure taxonomy, and one
technique-level finding, packaged as a single industrial-grade case study.

\begin{itemize}
  \item \textbf{C1 -- System.}
        We describe \agnes, a shooting-stage pipeline that composes
        \comfy\ graph execution, a locked style specification
        (\styledef), and face-embedding-based character anchoring, and
        the standards that make its outputs reproducible across operators.

  \item \textbf{C2 -- Failure Analysis at episode scale.}
        We release the audit records of \nts: \textbf{52 shot-level
        inspection records} covering six scenes and 18 fields each, of
        which 40 were rejected. We provide the first public distribution
        of rejection reasons in an AIGC animation pipeline, together with
        a taxonomy derived from them.

  \item \textbf{C3 -- Negative-prompt efficacy.}
        We quantify a counter-intuitive behaviour: constraint expressed
        as a \emph{negative} prompt is systematically \emph{less} likely
        to suppress the artefact it targets than the same constraint
        expressed positively. We report this as an ablation with a
        single manipulated variable.
\end{itemize}

\subsection{Why the failure distribution matters}

Two consequences follow directly from a 23\% one-pass rate. First, per-shot
pricing models systematically overstate deliverable output. Second, automatic
quality control cannot be designed from the success cases alone: using a
motion metric as an automatic reject rule, we observe that
\textbf{36 of 52} records are classified as \emph{motion-anomalous} — a
classification with no discriminative power. The dominant defects lie
elsewhere, and identifying them is the substance of \Cref{sec:failure}.

\subsection{Paper structure}

\Cref{sec:related} reviews prior work; \Cref{sec:system} describes the
pipeline; \Cref{sec:failure} presents the failure taxonomy;
\Cref{sec:ablation} reports the negative-prompt ablation;
\Cref{sec:discussion} discusses implications; \Cref{sec:conclusion}
concludes.

% TODO(写作顺序建议)
%   1. 先写本节 Motivation 段（已含 23% 这个数字，是全篇最有力的一句）
%   2. 再写 03_system（素材最全，写起来最顺）
%   3. 再写 04_failure（数据已实跑，见下方 FIXME）
%   4. 05_ablation 必须等实验跑完才能写完整
%   5. 02_related_work 最后写（它决定 C3 是否还成立）

% TODO
\section{Related Work}
% 待写。本文件是编译占位：缺少它 make 会因 \input 找不到文件而中止。

% TODO
\section{System}
% 待写。本文件是编译占位：缺少它 make 会因 \input 找不到文件而中止。

% TODO
\section{Failure Analysis}
% 待写。本文件是编译占位：缺少它 make 会因 \input 找不到文件而中止。

% TODO
\section{Ablation Study}
% 待写。本文件是编译占位：缺少它 make 会因 \input 找不到文件而中止。

% TODO
\section{Discussion}
% 待写。本文件是编译占位：缺少它 make 会因 \input 找不到文件而中止。

% TODO
\section{Conclusion}
% 待写。本文件是编译占位：缺少它 make 会因 \input 找不到文件而中止。


% ---------- MTAP 要求：材料可用性声明 ----------
\section{Declarations}
\begin{itemize}
  \item[--] \textbf{Funding}         % TODO  funding 或 "No funding received"
  \item[--] \textbf{Competing interests}
  \item[--] \textbf{Data and code availability}
        % 指向 github 仓库。未发表前只给 data/ code/ 路径，不给分析结论。
  \item[--] \textbf{Acknowledgments}
\end{itemize}

\bibliographystyle{splncs04}
\bibliography{bib/references}

\appendix
\section{Appendix A · Reproduction Details}
% 工作流 JSON 结构、审核指标定义、消融实验超参

\end{document}
